Industrial systems fail unpredictably, costing over \$50 billion annually in unplanned downtime~\cite{mckinsey2023}. Predictive Maintenance (PdM) leverages sensor data to forecast equipment failures before they occur, enabling proactive interventions that reduce costs and improve safety. Recent advances in Time Series Foundation Models (TSFMs)---large models pretrained on internet-scale temporal data---promise zero-shot transfer across domains without task-specific training. Models such as MOMENT~\cite{goswami2024moment}, Chronos~\cite{ansari2024chronos}, and TimeGPT~\cite{garza2023timegpt} have demonstrated impressive performance on standard forecasting benchmarks.

However, industrial PdM poses unique challenges that are absent from standard academic benchmarks:
\begin{itemize}
    \item \textbf{Non-IID distributions}: Sensor data varies significantly across machines, operating conditions, and deployment sites, violating the independent and identically distributed assumption that most TSFMs rely upon.
    \item \textbf{Privacy constraints}: Industrial sensor data often cannot leave factory premises due to proprietary and security concerns, limiting cloud-based model deployment.
    \item \textbf{Sparse failure labels}: Equipment failures are rare events, creating severe class imbalance that challenges supervised learning approaches.
    \item \textbf{Multi-variate sensor fusion}: Modern industrial systems generate data from dozens of correlated sensors that must be jointly modeled to capture system-level degradation patterns.
\end{itemize}

Despite these challenges, no systematic evaluation of TSFMs exists for industrial PdM scenarios. Existing benchmarks such as FoundTS~\cite{foundts2024} and GIFT-Eval~\cite{gifteval2024} focus on clean, academic datasets that fail to capture the realities of industrial deployment. This gap between benchmark performance and industrial applicability remains unquantified.

\subsection{Contributions}

This paper makes four key contributions:
\begin{enumerate}
    \item \textbf{First industrial PdM benchmark} of six state-of-the-art TSFMs across forecasting, anomaly detection, and Remaining Useful Life (RUL) prediction tasks on six public industrial datasets.
    \item \textbf{Novel SCADA-optimized preprocessing pipeline} that handles domain-specific gaps and improves cross-domain transfer performance by 18\%.
    \item \textbf{Taxonomy of TSFM industrial failures}: We document a 35\% MAE degradation on industrial data compared to benchmark performance, and a 62\% failure rate in zero-shot RUL prediction.
    \item \textbf{Federated readiness assessment} exposing privacy and edge deployment limitations that prevent current TSFMs from being deployed in real industrial settings.
\end{enumerate}

This work establishes the empirical foundation for developing industrially-robust TSFMs, revealing why current general-purpose models fail in real-world deployment scenarios and motivating the need for specialized approaches.
